\section{Semesta Pembicaraan (\textit{Universe of Discourse})}

Membiarkan variabel diisi oleh objek sembarang tanpa batas ruang lingkup dapat menimbulkan \textbf{ambiguitas} dan kesalahan interpretasi: predikat yang dimaksudkan untuk satu jenis entitas mungkin ``dipaksa'' pada entitas yang tidak relevan \cite{rosen2019}. Untuk itu diperkenalkan \textbf{semesta pembicaraan} (\textit{universe of discourse}).

\subsection{Definisi Semesta Pembicaraan}

\begin{definisi}[Semesta Pembicaraan / \textit{Universe of Discourse}]
\logic{Semesta pembicaraan} (disimbolkan $\mathcal{U}$) adalah himpunan objek yang \textbf{diperbolehkan} dan \textbf{dianggap relevan} sebagai nilai substitusi untuk variabel dalam fungsi proposisional pada konteks pembicaraan tertentu.
\end{definisi}

Penetapan $\mathcal{U}$ setara dengan menetapkan ``domain'' yang sah: hanya anggota $\mathcal{U}$ yang dijadikan kandidat konstanta pengganti variabel. Tanpa kesepakatan semesta, pihak pembaca dapat memilih substitusi yang tidak selaras dengan maksud penulis (misalnya memasukkan entitas non-numerik ke predikat aritmetika).

\subsection{Akibat Tanpa Batasan Semesta}

Perhatikan fungsi proposisional:
\[ P(x) \equiv \text{``}x \text{ adalah bilangan ganjil sehingga } x \bmod 2 = 1\text{''}. \]

\begin{enumerate}
    \item \textbf{Semesta } $\mathcal{U} = \mathbb{Z}^+$ \textbf{(bilangan bulat positif)}
    \begin{itemize}
        \item Contoh: $P(3)$ sesuai interpretasi bilangan ganjil $\implies$ \textbf{True}.
        \item Contoh: $P(6)$ tidak memenuhi ``ganjil'' $\implies$ \textbf{False}.
        \item Predikat tetap bermakna dalam kerangka bilangan.
    \end{itemize}

    \item \textbf{Substitusi di luar jenis yang dimaksud}
    \begin{itemize}
        \item Jika $x$ diganti ``Gunung Merapi'', kalimat yang terbentuk tidak lagi merupakan aplikasi predikat aritmetika yang wajar.
        \item Terjadi ketidakcocokan tipe atau makna (\textit{type mismatch} / kesalahan semantik) \cite{wiki_domain_of_discourse}.
    \end{itemize}
\end{enumerate}

\subsection{Pengaruh Semesta terhadap Nilai Kebenaran}

Pemilihan $\mathcal{U}$ dapat mengubah status kebenaran kalimat yang tampak sama secara sintaksis \cite{wiki_first_order_logic}.

\begin{contoh}
Misalkan $M$ menyatakan: untuk setiap $x$ yang dipertimbangkan, berlaku $x^2 \ge 0$.
\begin{itemize}
    \item Jika $\mathcal{U} = \mathbb{R}$, maka klaim tersebut \textbf{benar (True)} untuk semua $x \in \mathcal{U}$.
    \item Jika $\mathcal{U} = \mathbb{C}$, ambil $x = 3i$. Maka $x^2 = -9$, sehingga $x^2 \ge 0$ tidak terpenuhi dalam urutan bilangan riil; dalam interpretasi yang memakai bilangan kompleks, pernyataan serupa perlu dirumuskan ulang (misalnya dengan memperjelas domain dan makna ``$\ge$'').
\end{itemize}
\end{contoh}

Jadi, semesta pembicaraan bukan hanya catatan tambahan, melainkan bagian dari kesepakatan interpretasi yang menentukan apakah suatu formalisasi konsisten dan apakah evaluasi kebenarannya tepat.
